\subsection{REGULAR ELEMENTS FOR A LINEAR REPRESENTATION}

Lemma 1. Let M be an analytic manifold over k and \(a = (a_{0}, \ldots, a_{n-1}, a_{n} = 1)\) a sequence of analytic functions on M. For all \(x \in M\) , let \(r_{a}(x)\) be the upper bound of those \(i \in [0, n]\) such that \(a_{j}(x) = 0\) for j \textless{} i and let \(r_{a}^{0}(x)\) be the upper bound of those \(i \in [0, n]\) such that \(a_{j}\) is zero on a neighbourhood of x for j \textless{} i.

\begin{enumerate}
\def\labelenumi{(\roman{enumi})}
\item
  The function \(r_a\) is upper semi-continuous.
\item
  For all \(x \in \mathrm{M}\) , \(r_a^0(x) = \liminf_{y \to x} r_a(y)\) .
\item
  The function \(r_a^0\) is locally constant.
\item
  The set of points \(x \in M\) such that \(r_{a}^{0}(x) = r_{a}(x)\) is the set of points of M on a neighbourhood of which \(r_{a}\) is constant. This is a dense open subset of M. If k = C and M is finite dimensional and connected, it is open and connected.
\item
  If \(r_a(x) = i\) , then \(a_i(x) \neq 0\) and, for all \(y\) in a neighbourhood of \(x\) , we have \(a_i(y) \neq 0\) , so \(r_a(y) \leq i\) .
\item
  If \(r_{a}^{0}(x) = i\) , the functions \(a_{0}, \ldots, a_{i-1}\) are zero on a neighbourhood of x and, for any y in this neighbourhood, \(r_{a}(y) \geq i\) . Consequently, \(\liminf_{y \to x} r_{a}(y) \geq i\) . Every neighbourhood of x contains a point y such that \(a_{i}(y) \neq 0\) and hence \(r_{a}(y) \leq i\) . Thus \(\liminf_{y \to x} r_{a}(y) = i\) .
\item
  Let \(i = r_a^0 (x)\) and let \(V\) be a neighbourhood of \(x\) such that \(a_{j}(y) = 0\) for all \(y\in V\) and all \(j < i\) . Then \(x\in M - Z\) , where \(Z\) denotes the set of points of \(M\) in a neighbourhood of which the function \(a_{i}\) is zero. Since \(Z\) is closed in \(M\) (Differentiable and Analytic Manifolds, Results, 5.3.5), \(V\cap (M - Z)\) is a neighbourhood of \(x\) . For every point \(y\) in this neighbourhood, \(r_a^0 (y) = i\) .
\item
  The function \(r_{a}-r_{a}^{0}\) is upper semi-continuous and its value at any point is \(\geq0\) . If \(r_{a}(x)=r_{a}^{0}(x)\) , \(r_{a}-r_{a}^{0}\) is zero on a neighbourhood of x, which shows that \(r_{a}\) is constant on a neighbourhood of x by (iii). Conversely, if \(r_{a}\) is constant on a neighbourhood of x, then \(r_{a}^{0}(x)=r_{a}(x)\) by (ii). The set of points \(x\in M\) such that \(r_{a}^{0}(x)=r_{a}(x)\) is thus an open subset \(\Omega\) of M. If \(x\in M\) and if \(r_{a}^{0}(x)<r_{a}(x)\) , every neighbourhood of x contains a point y such that \(r_{a}(y)<r_{a}(x)\) and \(r_{a}^{0}(y)=r_{a}^{0}(x)\) . Every neighbourhood of x thus contains a point y such that
\end{enumerate}

\[
r _ {a} (y) - r _ {a} ^ {0} (y) <   r _ {a} (x) - r _ {a} ^ {0} (x).
\]

It follows that \(\Omega\) is dense in M.

If M is connected and if p is the value of \(r_{a}^{0}\) on M, the points of \(\Omega\) are the points \(x \in M\) such that \(a_{p}(x) \neq 0\) . If k = C, this implies that \(\Omega\) is connected by Lemma 3 of Appendix II.

Let \(\rho\) be an analytic linear representation of \(G\) on a vector space \(V\) of finite dimension \(n\) over \(k\) . Put

\[
\det (\mathrm{T} - \rho (g) + 1) = a _ {0} (g) + a _ {1} (g) \mathrm{T} + \dots + a _ {n - 1} (g) \mathrm{T} ^ {n - 1} + \mathrm{T} ^ {n}.
\]

The functions \(r_a\) and \(r_a^0\) associated to the sequence \((a_0, a_1, \ldots, a_{n-1}, 1)\) will be denoted by \(r_\rho\) and \(r_\rho^0\) , respectively. Then, for all \(g \in G\) ,

\[
\begin{array}{l} r _ {\rho} (g) = \dim \mathrm{V} ^ {1} (\rho (g)) \\ r _ {\rho} ^ {0} (g) = \liminf _ {g ^ {\prime} \to g} \dim \mathrm{V} ^ {1} (\rho (g ^ {\prime})). \end{array}
\]

Lemma 2. Let \(0 \to V' \to V \to V'' \to 0\) be an exact sequence of G-modules defined by analytic linear representations \(\rho', \rho, \rho''\) of G, respectively. Then:

\[
r _ {\rho} = r _ {\rho^ {\prime}} + r _ {\rho^ {\prime \prime}}, \text { and } r _ {\rho} ^ {0} = r _ {\rho^ {\prime}} ^ {0} + r _ {\rho^ {\prime \prime}} ^ {0}.
\]

Indeed, for all \(g \in \mathbf{G}\) , there is (§1, no. 1, Cor. 3 of Th.1) an exact sequence

\[
0 \to (\mathrm{V} ^ {\prime}) ^ {1} (\rho^ {\prime} (g)) \to \mathrm{V} ^ {1} (\rho (g)) \to (\mathrm{V} ^ {\prime \prime}) ^ {1} (\rho^ {\prime \prime} (g)) \to 0,
\]

which proves the first assertion. The second follows from it since, by Lemma 1 (iv), \(r_{\rho}^{0} = r_{\rho}, r_{\rho'}^{0} = r_{\rho'}\) and \(r_{\rho''}^{0} = r_{\rho''}\) on a dense open subset of G.

DEFINITION 1. An element \(g \in \mathbf{G}\) is called regular for the linear representation \(\rho\) if \(r_{\rho}(g) = r_{\rho}^{0}(g)\) .

PROPOSITION 1. The regular points for an analytic linear representation \(\rho\) of G are the points of G in a neighbourhood of which \(r_{\rho}\) is constant. They constitute a dense open subset of G. If k = C and G is connected, the set of regular points for \(\rho\) is connected.

This follows from Lemma 1 (iv).

Remark. Let \(G^{*}\) be an open subgroup of G. An element \(a \in G^{*}\) is a regular element of G for the linear representation \(\rho\) of G if and only if it is a regular element of \(G^{*}\) for the linear representation \(\rho|G^{*}\) .

